% !TEX root = ../notes.tex

\documentclass[../notes.tex]{subfiles}

\begin{document}

\section{April 15}
The homework has been posted.

\subsection{Comparison of Obstructions}
While we're here, we define some refinements of the descent obstruction.
\begin{notation}
	Fix a nice variety $X$ over a field $k$ and an algebraic group $G$. Then we define $X(\AA_k)^{\mathrm{\acute et},\mathrm{Br}}$ to be the intersection of
	\[\bigsqcup_{\sigma\in\mathrm H^1(k;G)}f^\sigma\left(Y^\sigma(\AA_k)^{\mathrm{Br}}\right)\]
	as $f$ varies over all $G$-torsors $f\colon Y\to X$ and finite \'etale group schemes $G$. We similarly define the  $X(\AA_k)^{\mathrm{\acute et},\mathrm{descent}}$.
\end{notation}
\begin{proposition} \label{prop:descent-is-et-br}
	Fix a nice variety $X$ over a global field $k$. Then
	\[X(\AA_k)^{\mathrm{descent}}=X(\AA_k)^{\mathrm{\acute et},\mathrm{descent}}=X(\AA_k)^{\mathrm{\acute et},\mathrm{Br}}\subseteq X(\AA_k)^{\mathrm{Br}}.\]
\end{proposition}
\begin{proof}
	Let's show inclusions in sequence.
	\begin{itemize}
		\item It turns out that the Brauer--Manin obstruction is a special case of the descent obstruction for the groups $\op{PGL}_n$. (This is basically \Cref{thm:gabber}.) This implies that
		\[X(\AA_k)^{\mathrm{descent}}\subseteq X(\AA_k)^{\mathrm{Br}}.\]
		Similarly, $X(\AA_k)^{\mathrm{\acute et},\mathrm{descent}}\subseteq X(\AA_k)^{\mathrm{\acute et},\mathrm{Br}}$.
		\item Using the structure theory of algebraic groups, one can show that $X(\AA_k)^{\mathrm{Br}}$ is the same as using obstructions from all connected groups. By further adding in some component group, one finds that
		\[X(\AA_k)^{\mathrm{\acute et},\mathrm{descent}}\subseteq X(\AA_k)^{\mathrm{\acute et},\mathrm{Br}}.\]
		\item It further turns out that $X(\AA_k)^{\mathrm{descent}}\subseteq X(\AA_k)^{\mathrm{\acute et},\mathrm{descent}}$, though this is not obvious. This is due to Skorobogatov and Stoll.
		\qedhere
	\end{itemize}
\end{proof}
\begin{remark}
	One could try to further define $X(\AA_k)^{\mathrm{descent},\mathrm{descent}}$, but it turns out to equal $X(\AA_k)^{\mathrm{descent}}$.
\end{remark}
\begin{remark}
	These obstructions are not sufficient.
	\begin{itemize}
		\item Skorobagatov found a nice variety $X$ where $X(\AA_k)^{\mathrm{Br}}$ is nonempty while $X(\AA_k)^{\mathrm{descent}}$ is empty. The example is found among bi-elliptic surfaces.
		\item Bjorn Poonen showed that $X(\AA_k)^{\mathrm{descent}}$ may be nonempty when $X(k)$ is empty. His original counterexample was three-dimensional, but there may be two-dimensional counterexamples.
	\end{itemize}
\end{remark}
\begin{remark}
	Bjorn Poonen has a heuristic that $X(\AA_k)^{\mathrm{Br}}$ is sufficient to detect whether curves have rational points, which explains why these counterexamples are found among surfaces and three-folds.
\end{remark}
\begin{remark}
	In class, we discussed what might happen if we try to add in abelian varieties. It is not obvious if this refines the descent obstruction.
\end{remark}

\subsection{Descent Is Not Enough}
For the rest of class, we will focus on giving an example of some nice $X$ where $X(\AA_k)^{\mathrm{descent}}$. Poonen's original example was a bundle of Ch\^atelet surfaces, but we will present an easier example (due to Coliot-Th\'el\`ene, P\'al, and Skorobogatov) of a bundle of quadrics.
\begin{theorem}
	There is a nice variety $X$ over $\QQ$ for which $X(k)=\emp$ while $X(\AA_k)^{\mathrm{descent}}\ne\emp$.
\end{theorem}
\begin{proof}
	We proceed in steps. Let's start with the construction.
	\begin{enumerate}
		\item Start with a nice $\QQ$-curve $C$ such that $C(\QQ)$ is a single point $c$. For example, a smooth projective model of $y^2=x^3-3$ will work, which one shows by doing a $2$-Selmer group calculation to show rank zero and then computing the torsion group by computing $C(\FF_p)$ for various primes $p$.

		The smoothness (and properness) of $C$ allows us to find a map $f\colon C\to\PP^1$ which is \'etale at $c$, and we may as well assume that $f(c)=\infty$. For $C\colon y^2=x^3-3$, the map $y\colon C\to\PP^1$ will do. While we're here, we also choose a simply connected open neighborhood $U\subseteq C(\RR)$ of $c$, and we choose it small enough so that it remains simply connected and open as $f(U)\subseteq\PP^1_\RR$. By changing coordinate again, we may assume that $1\in f(U)$ and that $f$ is \'etale above $0$ and $1$.

		\item Now, glue the two bundles $Y^{(t)}\subseteq\PP^4\times\AA^1_t$ and $Y^{(T)}\subseteq\PP^4\times\AA^1_T$ given by
		\[\begin{cases}
			Y^{(t)}\colon t(t-1)x_0^2+\left(x_1^2+x_2^2+x_3^2+x_4^2\right)=0, \\
			Y^{(T)}\colon (1-T)X_0^2+\left(x_1^2+x_2^2+x_3^2+x_4^2\right)=0.
		\end{cases}\]
		Namely, these may be glued along the equations $t=1/T$ and $x_0=X_0T$, which gives a quadric bundle $Y\to\PP^1_{t}$. This has two degenerate fibers, which are above $0$ and $1$; namely, we get some degenerate cones above those points.

		\item We construct our $X$. By pulling back $Y\to\PP^1$ along $f\colon C\to\PP^1$, we obtain a quadric bundle $\pi\colon X\to C$. Note that $Y\to\PP^1$ is certainly projective with geometrically irreducible fibers, so we receive the same for $X\to C$. In fact, $Y\to\PP^1$ is smooth away from the fibers of $0$ and $1$, and $C\to\PP^1$ is smooth above $0$ and $1$, so $X$ is smooth over $\PP^1$ and thus over $\QQ$.
	\end{enumerate}
	We now run our checks on $X$.
	\begin{enumerate}[resume]
		\item Let's show that $X(\QQ)=\emp$. On one hand, any $x\in X(\QQ)$ provides a $\QQ$-point of $C$, so $\pi(x)=c$. But then $f(x)=\infty$. On the other hand, $x$ must also map to a rational point of $Y$, which now must be in the fiber $Y_\infty\subseteq Y$. But taking $T=0$ in $Y^{(T)}$ finds no real points, so there are no rational points here.

		\item We now want to show that there are local points which survive the descent obstruction. Let's start by showing that there are local points.
		\begin{itemize}
			\item We handle finite places. To begin, we recall that every quadratic form over $\QQ_p$ in at least five variables has a nontrivial zero.\footnote{This is false for four variables because one can use a quadratic form in four variables to cut out a quaternion algebra.} Indeed, this is classical. It is basically Hensel's lemma. For example, avoiding $p=2$ (which is handled separately), we may diagonalize and reorganize the coefficients, so we may assume that the quadratic form looks like $u_1x_1^2+u_2x_2^2+u_3x_3^2+u_4x_4^2+u_5x^5$, where $u_1$, $u_2$, and $u_3$ are units. Then a counting argument finds a smooth point over $\FF_p$, which then lifts to $\ZZ_p$.

			Now, we define $x_p\coloneqq(y_p,c)\in X(\QQ_p)$, where $y_p$ is a distinguished point in $y_p\in Y_\infty(\QQ_p)$.

			\item At the infinite place, we let $y_\RR$ be the unique real point of $Y_1(\RR)$, and we choose $c_\RR\in U$ for which $f(c_\RR)=1$ in $\PP^1(\RR)$. Thus, we may define $x_\RR\coloneqq(y_\RR,c_\RR)$ to be a point in $X(\RR)$.
		\end{itemize}
		For future use, we set $x\coloneqq(x_\RR,x_2,x_3,\ldots)$ to be this point of $X(\AA_\QQ)$.

		\item As a warm-up, we claim that $x\in X(\AA_\QQ)^{\mathrm{Br}}$. Well, select any $A\in\op{Br}X$, and we want to show that $x\in X(\AA_\QQ)^A$. For this, recall from \Cref{prop:br-quadric-bundle} that $\op{Br}C\to\op{Br}X$ is surjective, so we may find $A_C\in\op{Br}C$ which maps to $A$. Thus,
		\[\sum_v\op{inv}_v\op{ev}_A(x_v)=\sum_v\op{inv}_v\op{ev}_{A_C}(\pi(x_v))\]
		by unwinding the evaluations. Now, for finite $v$, we see that $\pi(x_v)=c$. Even at $\infty$, one knows that $\op{ev}_{A_C}$ is some continuous map $C(\RR)\to\frac12\ZZ/\ZZ$, so $\op{ev}_{A_C}(c_\RR)=\op{ev}_{A_C}(c)$ still. Thus, this sum is
		\[\sum_v\op{inv}_v\op{ev}_{A_C}(c),\]
		which vanishes because $c$ is actually a rational point on $C$!

		\item We now claim that $x\in X(\AA_\QQ)^{\mathrm{\acute et},\mathrm{Br}}$, which will complete the proof by \Cref{prop:descent-is-et-br}. Note that the fibers of $X\to C$ are all simply connected (they are quadric hypersurfaces), so every \'etale over of $X$ comes from an \'etale cover of $C$. Now, choose a finite \'etale group scheme $G$, and choose a $G$-torsor $X'\to X$. Because $G$ is \'etale, the previous few sentences explain that this torsor must come from a $G$-torsor $C'\to C$. The descent obstruction does not care if we replace $X'$ with a twist, so we may twist $C'$ so that $c\in C(\QQ)$ lifts to some $c''\in C'$; accordingly, we let $C''$ be the irreducible component of $C'$.\footnote{There is a classification of the \'etale covers of elliptic curves, which tells us that $C''$ is an elliptic curve with identity $c''$, and the map $C''\to C$ is an isogeny.} All of this fits into the following pullback squares.
		% https://q.uiver.app/#q=WzAsOCxbMCwwLCJYJyciXSxbMCwxLCJDJyciXSxbMSwwLCJYJyJdLFsxLDEsIkMnIl0sWzIsMCwiWCJdLFsyLDEsIkMiXSxbMywwLCJZIl0sWzMsMSwiXFxQUF4xIl0sWzMsNSwiRyJdLFsyLDQsIkciXSxbNiw3XSxbNCw2XSxbNSw3XSxbNCw1XSxbMiwzXSxbMSwzXSxbMCwxXSxbMCwyXV0=&macro_url=https%3A%2F%2Fraw.githubusercontent.com%2FdFoiler%2Fnotes%2Fmaster%2Fnir.tex
		\[\begin{tikzcd}[cramped]
			{X''} & {X'} & X & Y \\
			{C''} & {C'} & C & {\PP^1}
			\arrow[from=1-1, to=1-2]
			\arrow[from=1-1, to=2-1]
			\arrow["G", from=1-2, to=1-3]
			\arrow[from=1-2, to=2-2]
			\arrow[from=1-3, to=1-4]
			\arrow[from=1-3, to=2-3]
			\arrow[from=1-4, to=2-4]
			\arrow[from=2-1, to=2-2]
			\arrow["G", from=2-2, to=2-3]
			\arrow[from=2-3, to=2-4]
		\end{tikzcd}\]
		Now, $C''$ is smooth and projective, and it has a rational point, so it follows that $C''$ is a nice curve. (In particular, it is geometrically irreducible.) By pulling back along $X\to C$, we see that $X''$ is also nice.

		It is now enough to show that $x\in X(\AA_k)$ lifts to some $x''\in X''(\AA_k)^{\mathrm{Br}}$. To start, let's show that it lifts to $X''(\AA_k)$.
		\begin{itemize}
			\item For finite prime $p$, we define $x''_p\coloneqq(x_p,c'')$.
			\item At $\infty$, we note that $U$ is simply connected, so its inverse image in $C''(\RR)$ must be a disjoint union of copies of $U$; accordingly, we let $c''_\RR$ be the pre-image of $c_\RR\in U$ in the same copy of $U$ in $C''(\RR)$ as $c''$.
		\end{itemize}
		These local points now glue to $x''\in X''(\AA_k)$, so it remains to check that $x''$ is in $X'(\AA_k)^{\mathrm{Br}}$. Because $X''$ is just a connected component of $X'$, new Azumaya algebras gained from going from $X''$ to $X'$ will not provide a meaningful obstruction, so it is enough to check that $x''\in X''(\AA_k)^{\mathrm{Br}}$. Well, one can just repeat the same argument from the previous step word-for-word, so we are done!
		\qedhere
	\end{enumerate}
\end{proof}
\begin{remark}
	One could try to perturb this example by replacing $C$ with a surface, but nobody knows an example of a simply connected surface with exactly one rational point.
\end{remark}
We end the course with a conjecture.
\begin{conj}[Colliot-Th\'el\`ene, Karevsky, Sansuc] \label{conj:br-enough}
	Fix a nice variety $X$ over a number field which is geometrically rationally connected. If $X(\AA_k)^{\mathrm{Br}}\ne\emp$, then $X(k)\ne\emp$.
\end{conj}
\begin{remark}
	Rationally connected means that any two points can be connected by some map $\PP^1\to X$.
\end{remark}
\begin{remark}
	Basically nothing is known about \Cref{conj:br-enough}. It is known for Ch\^atelet surfaces.
\end{remark}
\begin{remark}
	One expects that $\op{Br}X=\op{Br}k$ for varieties of general type, so $X(\AA_k)^{\mathrm{Br}}=X(\AA_k)$. On the other hand, it is rather rare for such varieties to admit rational points.
\end{remark}

\end{document}